\documentclass[11pt]{article}
\usepackage[a4paper,margin=21mm]{geometry}
\usepackage[no-math]{fontspec}
\setmainfont{Latin Modern Roman}
\newfontfamily\CJKfont{Noto Serif CJK SC}
\XeTeXlinebreaklocale "zh"
\XeTeXlinebreakskip=0pt plus 1pt
\usepackage{amsmath,amssymb,amsthm,amscd,graphicx,color}
\usepackage{xurl}
\usepackage[unicode,hidelinks]{hyperref}
\allowdisplaybreaks
\setlength\emergencystretch{3em}
\numberwithin{equation}{section}

\numberwithin{equation}{section}
\newtheorem{Theorem}{Theorem}[section]
\newtheorem*{Theorem*}{Theorem}
\newtheorem{Corollary}[Theorem]{Corollary}
\newtheorem{Lemma}[Theorem]{Lemma}
\newtheorem{Proposition}[Theorem]{Proposition}
 { \theoremstyle{definition}
\newtheorem{Definition}[Theorem]{Definition}
\newtheorem{Note}[Theorem]{Note}
\newtheorem{Example}[Theorem]{Example}
\newtheorem{Remark}[Theorem]{Remark} }

\begin{document}
\begin{center}\Large Torus restrictions of left-invariant SKT metrics on compact semisimple Samelson complex Lie groups\end{center}
\noindent\textbf{Stable ID:} 1092\quad\textbf{Author:} Anna Fino; Gueo Grantcharov\par
\noindent\textbf{Work:} CYT and SKT Metrics on Compact Semi-Simple Lie Groups\par
\noindent\textbf{Source:} \url{https://sigma-journal.com/2023/028/}\par
\noindent\textbf{Version:} Official SIGMA 19 (2023) article 028, DOI 10.3842/SIGMA.2023.028; journal PDF and native TeX acquired 2026-09-30, exact bytes pinned by SHA256\par
\noindent\textbf{Rights:} CC BY 4.0. Authors retain copyright. \url{https://creativecommons.org/licenses/by/4.0/}. Reuse and typesetting adaptation; original language and mathematical content preserved.\par
\noindent\textbf{Difficulty (prior selection preserved):} {\CJKfont H2: The proof turns the SKT equation into a quadratic relation among the curvature classes of the maximal-torus fibration. It separates the simple factors through the flag-manifold cohomology relations and compares the metric change-of-basis matrix with the factorwise Killing forms. This produces both the exact torus restriction and the compatibility obstruction for the Samelson complex structure, beyond constructing a standard invariant example.}\par
\medskip\noindent\textbf{Editorial boundary note (not author text):} Complete original Theorem3.2 and its full torus-restriction/complex-structure obstruction proof. The Hermitian/SKT/fundamental-form conventions, full compact semisimple/Samelson root and Killing-form normalizations, simple-factor decomposition, all metric and maximal-torus hypotheses, current symmetrization proof, complete curvature-class relation and full transition-matrix comparison are retained. The precise author-cited torus-bundle curvature identity and standard Chevalley flag-cohomology relation remain stated author prerequisites; no new proof is substituted. The author expressly announces this as the first global SKT-compatibility consequence, a standalone obstruction outcome rather than an adjacent technical lemma counted with the later CYT/SKT classification. The later characterization corollary, SO(9) example and CYT/Bismut-flat classification are excluded, with supporting symmetrization receiving no extra count. Literal source cohomology degree and matrix indices are preserved without mathematical adjudication.\par
\medskip\hrule\medskip\noindent\textbf{Author text begins below.}\par
\setcounter{section}{1}
% Original source lines 68--78
On every Hermitian manifold $(M, I, g)$ there exists a one-parameter family of Hermitian connections which can be distinguished by their torsion tensor and coincide with the Levi-Civita connection when $M$ is K\"ahler. Among them there are
the Chern connection $\nabla^{\rm C}$ on the holomorphic tangent bundle and the Bismut (or Strominger) connection $\nabla^{\rm B}$, which is also called by physicists the K\"ahler with torsion (KT) connection. $\nabla^{\rm B}$ is the unique Hermitian connection whose torsion tensor is totally skew-symmetric and its torsion $T^{\rm B}$ is characterized by the condition
\[
g \big(T^{\rm B}(X,Y), Z\big) = {\rm d}F(IX,IY,IZ),
\]
where $F$ is the fundamental form $F(X,Y) = g(IX,Y)$, and $X$, $Y$, $Z$ are smooth vector fields. Since $\nabla^{\rm B}$ is a Hermitian connection, its holonomy is contained in the unitary group~${\rm U}(n)$, where~$n$ is the complex dimension of $M$. If the holonomy of $\nabla^{\rm B}$ is reduced to ${\rm SU}(n)$, the metric~$g$
is said to be Calabi--Yau with torsion (shortly CYT). It is also called {\it Bismut Ricci-flat}.
This type of geometry in the physical context was considered first by Strominger~\cite{Strominger} and
Hull~\cite{Hull}. In~\cite{GIP}, it was conjectured that every compact complex manifold with vanishing first Chern class admits a~CYT metric. Counterexamples to this conjecture appear in~\cite{FG}. In~\cite{Grantcharov} it has been shown that every compact complex homogeneous space with vanishing first Chern class, after an appropriate deformation of the complex structure, admits a homogeneous CYT metric, provided that the complex homogeneous space also has an invariant volume form.

{\em SKT} metrics, which are called also {\em pluriclosed}, are Hermitian metrics whose fundamental form $F$ is $\partial \overline \partial$-closed or equivalently, whose torsion $3$-form of the Bismut connection is closed. Interest in SKT manifolds stems from various sources~\cite{Swann}. They occur in physics in the context of supersymmetric theories, see for example \cite{GHR,HP,Strominger}. Secondly, every conformal class of a~hermitian metric on a~compact complex surface contains an SKT metric \cite{Gau}, but this property is no longer true in higher dimensions. Recently the SKT and CYT metrics have been studied in relation to the pluriclosed flow and the generalized Ricci flow \cite{AFSU, Barbaro, GJS, GS}.

\setcounter{section}{1}
% Original source lines 92--210
\section{Preliminaries on compact semi-simple Lie groups}
It was shown independently by Samelson \cite{Samelson} and Wang \cite{Wang} that every compact Lie group $G$ of even dimension admits left-invariant complex structures $I$.
Moreover, there exists a maximal toral subgroup $T$ of $G$ that is $I$-complex and the quotient $G/T$
is a complete flag variety of $G$ with a complex structure such that the quotient map $\pi\colon G \rightarrow G/T$ is holomorphic.

We now review the Samelson's construction of left-invariant complex structures on compact semi-simple Lie groups.
We will follow the conventions of Alekseevsky--Perelomov \cite{AlP}, and Bordemann--Forger--Romer \cite{BFR}.
 Let $G^{\rm c}$ be the complexification of a compact semi-simple Lie group $G$. Then $G$ may be regarded as a Lie subgroup of $G^{\rm c}$, and the Lie algebra of $G^{\rm c}$ is the complexification $\mathfrak g^{\rm c}$ of the Lie algebra $\mathfrak g$ of $G$.

If we fix a maximal abelian subalgebra $\mathfrak t$ of $\mathfrak g$, then $\mathfrak h= \mathfrak t^{\rm c}$ is a maximal abelian subalgebra of~$\mathfrak g^{\rm c}$ and we have the root space decomposition
\[
\mathfrak g^{\rm c} = \mathfrak{h} \oplus \bigoplus_{\alpha \in \Phi} \mathfrak g_{\alpha},
\]
where $\Phi$ is the finite subset of nonzero elements of ${\mathfrak h}^*$ called roots, and each $\mathfrak g_{\alpha}$ is the complex $1$-dimensional subspace of $\mathfrak g^{\rm c} $ given by
\[
\mathfrak g_{\alpha} := \{ X \in {\mathfrak g}^{\rm c} \mid [H , X] = \alpha (H) X, \, \forall H \in \mathfrak{h} \}.
\]

Therefore, there exists a system of positive roots, that is a set $P \subset \Phi$ satisfying $P \cap (- P) = \varnothing$, $ P \cup (-P)=\Phi$, and $\alpha, \beta \in P$, $ \alpha + \beta \in \Phi$ $\Rightarrow$ $\alpha + \beta \in P$, and such that ${\rm {span}}_{\mathbb C}(P) = \mathfrak h$. We note that $\Phi$ is also invariant under a set of reflections which leads to the classification of all irreducible root systems. Moreover, the structure of $\Phi$ determines the structure of $\mathfrak g^{\rm c}$ in a sense that $[\mathfrak g_{\alpha}, \mathfrak g_{\beta}] \subseteq \mathfrak g_{\alpha +\beta}$ when $\alpha+\beta \in \Phi$, $[\mathfrak h,\mathfrak g_{\alpha}] = \mathfrak g_{\alpha}$ and $[\mathfrak g_{\alpha}, \mathfrak g_{-\alpha}] \subseteq \mathfrak h$ and the other commutators vanish.

Among the roots in $P$ there is a set of simple roots, which we denote by $\alpha_1, \alpha_2,\dots ,\alpha_r$ and forms a basis of the space $\mathfrak t^*$. For any $\alpha \in P$, we have $\alpha = \sum_{i =1}^r m_i\alpha_i$, with $m_i$ being a non-negative integer. The basis $(\alpha_l)$ also determines a partial order of the set of positive roots and~$r$ is called rank of $\mathfrak g^{\rm c}$.

Denote by $B(X,Y) = {\rm tr} ({\rm ad}_X \circ {\rm ad}_Y)$ the Killing--Cartan bilinear form of $\mathfrak g^{\rm c}$. The restriction of $B$ to $\mathfrak t$ is non-degenerate and $\mathfrak g_{\alpha}$ is $B$-orthogonal to $\mathfrak g_{\beta}$ for $\beta \neq - \alpha$.
 For every $\alpha \in P$ we choose $(E_{\pm \alpha}, H_{\alpha})$ such that $\mathfrak g_{\pm \alpha}= {\rm {span}}_{\mathbb C} \langle E_{\pm\alpha}\rangle$, $[E_{\alpha}, E_{-\alpha}]=H_{\alpha}$, $B(E_{\alpha}, E_{-\alpha}) = 1$, $[H, E_{\alpha}] = \alpha(H)E_{\alpha}$ and $B(H, H_{\alpha}) = \alpha(H)$ for all $H \in \mathfrak t$.

Note that $\mathfrak t = {\rm span}_{\mathbb R} \langle {\rm i} H_{\alpha}\rangle$ is the Lie algebra of $T$. Moreover, since all compact forms of a~complex semi-simple Lie algebra are conjugate, we can assume that $\mathfrak g$ is the real span of $X_{\alpha}$,~$Y_{\alpha}$,~${\rm i}H_{\alpha}$, where $X_{\alpha} = \frac{1}{\sqrt{2}}(E_{\alpha}+E_{-\alpha}), Y_{\alpha} = \frac{1}{\sqrt{2}{\rm i}}(E_{\alpha}-E_{-\alpha}).$
In particular, $\{ X_{\alpha}, Y_{\alpha} \} $ is an orthonormal set with respect to $- B$ on $G$.


When the rank $r$ is even, Samelson proved that a choice of $P$ and of a complex structure on~$\mathfrak t $ determine a left-invariant complex structure $I$ on $G$, which is given by $I(E_{\alpha})= {\rm i}E_{\alpha}$, $I(E_{-\alpha})=-{\rm i}E_{-\alpha}$, for every $\alpha \in P$, and by a linear endomorphism of $\mathfrak{t}$ with square $-{{\rm Id}}$.

Note that if $ \mathfrak g = \sum_{i =1}^m \mathfrak g_i$, where $\mathfrak g_i$ are simple, then $\mathfrak t= \sum_{i=1}^m \mathfrak t_i$ and $B = \sum_{i=1}^m B_i$. We denote by $E_{\alpha}^{(i)}$, $E_{-\alpha}^{(i)}$ the corresponding root vectors in $\mathfrak g_i$ with respect to $\mathfrak t_i$ and by $P_i$ the set of positive roots of the simple factor $G_i$ with respect to $\mathfrak t_i$. We will not use the upper index when there is no confusion.

\begin{Remark} \label{remarkAkekseevski} Let $I$ be a Samelson complex structure on $G$. If $h$ is a left-invariant $I$-Hermitian metric on $G$ which is also invariant by the right $T$-action on $G$, then there exists a metric on $G/T$ such that $\pi\colon G\rightarrow G/T$ is a Riemannian submersion. Since the spaces $\langle X_{\alpha}, Y_{\alpha}\rangle$ are $T$-invariant and irreducible,
the metric $h$ can be written, up to a constant, as $h(X,Y)= -B (\Lambda(X),Y)$, where~$\Lambda$ is
a symmetric operator defined by $\Lambda(E_{\alpha})=\lambda_{\alpha} E_{\alpha}$ for every $\alpha \in P$ and any
appropriate matrix on $ \mathfrak t$ such that $h$ is positive definite (see, e.g., \cite{Arv}).
\end{Remark}

\section{Invariant SKT metrics}

 We can apply the symmetrization with respect to the right $T$-action to prove that if we have a~left-invariant SKT $I$-Hermitian metric,
then it is possible to construct a SKT metric which is right $T$-invariant.

\begin{Lemma}\label{llemma1}
Let $G$ be a compact semi-simple Lie group of even rank endowed with a Samelson complex structure $I$. Let $g$ be a left-invariant SKT $I$-Hermitian metric.
Then symmetrizing~$g$ by the right $T$-action on $G$, we obtain a Hermitian metric $h$ which is still SKT and also left-invariant and right $T$-invariant.
\end{Lemma}


\begin{proof} Since $I$ is invariant under the right $T$-action by the aforementioned result of Wang \cite{Wang}, we can symmetrize $g$. Indeed, let $d \nu$ be the standard volume on the torus $T$ and denote by $R_h$ the right translation by $h \in T$ on $G$. We can define
\[
 \overline{g}_x ({\bf u}, {\bf v}) := \int_T (R_h^* g)_x ({\bf u}, {\bf v}) \, {\rm d} \nu (h),
\]
for every $x \in G, {\bf u}, {\bf v} \in T_x G$. By construction $\overline{g}$ is invariant by the right $T$-action on $G$ and compatible with $I$. Moreover, since left and right translations commute and right $T$- translations are holomorphic, $\overline{g}$ is left-invariant.

Since the SKT condition ${\rm d}{\rm d}^{\rm c} \omega= 0$ is linear in $\omega$ and ${\rm d}$, ${\rm d}^{\rm c}$ commute with $R_h^*$, for every $ h \in T$, $\overline{g}$ is SKT. \end{proof}

Given a Samelson complex structure $I$, we can describe all possible left-invariant SKT $I$-Hermitian metrics which are also invariant by the right $T$-action on $G$.

\begin{Theorem}\label{lemma2}
Let $G$ be a compact semi-simple Lie group of even rank endowed with a Samelson complex structure $I$ and let $g$ be a left-invariant SKT $I$-Hermitian metric.
Then the restriction of $g$ on the fibers of $\pi\colon G \rightarrow G/T$ coincides with the restriction of a bi-invariant metric on $G$. As a consequence, the complex structure $I$ has to be compatible with a bi-invariant metric $g_0$.
\end{Theorem}

\begin{proof} By Lemma \ref{llemma1}, using the symmetrization it is not restrictive to suppose that $g$ is invariant under the right $T$-action.
 Suppose that $G$ up to a finite cover is given by the product $G_1\times G_2\times \dots \times G_m$ of simple Lie groups $G_i$ ($ \dim G_i > 1$). Then $\mathfrak t = \sum_{i = 1}^m \mathfrak t_i$, with $\mathfrak t_i$ the Cartan subalgebra of the Lie algebra $\mathfrak g_i$ of every factor $G_i$.


Select a $g$-unitary basis $(\theta_1, \dots , \theta_{2k})$ of $\mathfrak t^*$, such that $ I(\theta_{2l-1})=\theta_{2l}$.
Then by \cite{GGP}, we have $F = \sum_{l = 1}^k \theta_{2l-1}\wedge\theta_{2l} + \pi^*(F_b),$ where $F_b$ is a $(1,1)$-form of an invariant Hermitian structure on~$G/T$. Also by \cite[Section~5]{GGP},
\[
{\rm d}{\rm d}^{\rm c} F = \pi^*{\rm d}{\rm d}^{\rm c} F_b - \sum_{l = 1}^k \big(\omega_{2l-1}^2+\omega_{2l}^2\big),
\]
 where $\omega_i = {\rm d}\theta_i$ are the curvature forms of the connection on $\pi\colon G\rightarrow G/T$ determined by $\theta_i$. This in particular means that
 \[
 \sum_{l=1}^k \big([\omega_{2l-1}]^2+[\omega_{2l}]^2\big) = \sum_{i=1}^{2k} [{\rm d}\theta_i]^2 = 0
 \]
 in the cohomology $H^2(G/T, \mathbb{R})$.

Now, the flag manifold $G/T$ decomposes as $G/T = G_1/T_1\times\dots \times G_m/T_m$ and by the results of Chevalley \cite{Chevalley}, $H^2(G_i/T_i, \mathbb{R})$ is generated by the (exterior derivatives of the) simple roots in~$\mathfrak t_i^*$, and moreover, there is a unique quadratic relation among them. This quadratic relation can be written as $\sum_{j=1}^{s_i} \big[{\rm d}\widetilde{\theta}^{(i)}_j\big]^2 = 0$ for
a $-B_i$-orthonormal basis $\big(\widetilde{\theta}^{(i)}_1, \dots, \widetilde{\theta}^{(i)}_{s_i}\big)$ of $\mathfrak t_i^*$, which can be easily checked by direct calculations. Since $B =\sum_{i =1}^m B_i$, the basis is also $-B$-orthonormal.

It follows that the only quadratic relations among the cohomology classes in $H^2(G/T, \mathbb{R})$ are given by
\begin{equation}\label{quadraticrel}
 \sum_{i=1}^m c_i\sum_{j=1}^{s_i} \big[{\rm d}\widetilde{\theta}^{(i)}_j\big]^2 = 0, \end{equation} for some constants $c_1, \dots, c_m$.

Then
\begin{equation} \label{quadrelation}
\sum_{i=1}^{2k} [{\rm d}\theta_i]^2 = \sum_{i=1}^m c_i\sum_{j=1}^{s_i} \big[{\rm d}\widetilde{\theta}^{(i)}_j\big]^2.
\end{equation}

Denote the $-B$-orthonormal basis $\big(\widetilde{\theta}_1^{(1)}, \dots, \widetilde{\theta}_{s_1}^{(1)}, \dots, \widetilde{\theta}_1^{(m)}, \dots, \widetilde{\theta}_{s_m}^{(m)}\big)$ by $\sigma_l$ for $l=1,2,\dots ,2k$. Then $\theta_i = \sum_{l=1}^{2k} q_{li}\sigma_l$ for a transition matrix $Q = (q_{ij})$
and we can rewrite $\sum_{i=1}^{2k} [{\rm d}\theta_i]^2$ as
\[
\sum_{i=1}^{2k} [{\rm d}\theta_i]^2 = \sum_{i = 1}^{2k} \sum_{l,j =1}^{2k} q_{li} q_{ji} [{\rm d}\sigma_l \wedge {\rm d} \sigma_j].
\]
 On the other hand,
\[
 \sum_{i=1}^m c_i\sum_{j=1}^{s_i} \big[{\rm d}\widetilde{\theta}^{(i)}_j\big]^2 = c_1 \sum_{i =1}^{s_1} [{\rm d} \sigma_i]^2 + c_2 \sum_{i =1}^{s_2} [{\rm d} \sigma_{s_1 +i}]^2 + \dots + c_m \sum_{i =1}^{s_m} [{\rm d} \sigma_{s_1 + \dots + s_{m-1} + i} ]^2.
\]
Therefore, from the equality \eqref{quadrelation} we have
\[
 \sum_{i = 1}^{2k} \sum_{l,j =1}^{2k} q_{li} q_{ji} [{\rm d} \sigma_l \wedge {\rm d} \sigma_j] = c_1 \sum_{i =1}^{s_1} [{\rm d} \sigma_i]^2 + c_2 \sum_{i =1}^{s_2} [{\rm d} \sigma_{s_1 +i}]^2 + \dots + c_m \sum_{i =1}^{s_m} [{\rm d} \sigma_{s_1 + \dots + s_{m-1} + i} ]^2.
\]
Using that $[{\rm d}\sigma_l\wedge {\rm d}\sigma_j]$ are independent, for every $l\neq j$ we get that
\[
\sum_{i = 1}^{2k} q_{li} q_{ji} =0, \qquad l \neq j,
\]
and by \eqref{quadraticrel}
\[
\sum_{i =1}^{s_1} (q_{li})^2 = c_1, \quad \dots, \quad \sum_{i =s_{m-1}}^{s_m} (q_{li})^2 = c_m,
\]
i.e., that $QQ^t= D$, where $D$ is a diagonal matrix with entries $c_i$. From here we conclude, that~$c_i > 0$ and $Q$ preserves the metric $-\sum_{i = 1}^m c_i B_i$. In particular $\theta_i$ are $-\sum_{i=1}^m c_i B_i$-or\-thonor\-mal. As the $\theta_i$s are $g$- and $-\sum_{i=1}^m c_i B_i$-orthonormal, the restriction of $g$ to $\mathfrak t$ coincides with the restriction of a biinvariant metric, the theorem follows.
\end{proof}


\begin{Remark} It is well known that on an even dimensional compact Lie group equipped with a~bi-invariant metric $g_0$ there is always a $g_0$-compatible complex structure, which is necessary SKT and also Bismut-flat. But there exist Samelson's complex structures which are not compatible with any bi-invariant metric on $G$ as long as the rank of $G$ is at least 2.
\end{Remark}
\begin{thebibliography}{99}
\bibitem[1]{AlP}
Alekseevskii D.V., Perelomov A.M., Invariant {K}\"ahler--{E}instein metrics on
 compact homogeneous spaces, \href{https://doi.org/10.1007/BF01078469}{\textit{Funct. Anal. Appl.}} \textbf{20} (1986),
 171--182.

\bibitem[2]{AFSU}
Apostolov V., Fu X., Streets J., Ustinovskiy Yu., The generalized K\"ahler
 Calabi--Yau problem, \href{https://arxiv.org/abs/2211.09104}{arXiv:2211.09104}.

\bibitem[3]{Arv}
Arvanitoyeorgos A., An introduction to {L}ie groups and the geometry of
 homogeneous spaces, \textit{Stud. Math. Libr.}, Vol.~22, \href{https://doi.org/10.1090/stml/022}{Amer. Math. Soc.},
 Providence, RI, 2003.

\bibitem[4]{Barbaro}
Barbaro G., Global stability of the {P}luriclosed flow on compact
 simply-connected simple {L}ie groups of rank two, \href{https://doi.org/10.1007/s00031-022-09761-5}{\textit{Transform. Groups}}
 (2022),, \href{https://arxiv.org/abs/2202.13199}{arXiv:2202.13199}.

\bibitem[5]{BFR}
Bordemann M., Forger M., R\"omer H., Homogeneous {K}\"ahler manifolds: paving
 the way towards new supersymmetric sigma models, \href{https://doi.org/10.1007/BF01221650}{\textit{Comm. Math. Phys.}}
 \textbf{102} (1986), 605--647.

\bibitem[6]{Chevalley}
Chevalley C., Invariants of finite groups generated by reflections,
 \href{https://doi.org/10.2307/2372597}{\textit{Amer.~J.~Math.}} \textbf{77} (1955), 778--782.

\bibitem[7]{FG}
Fino A., Grantcharov G., Properties of manifolds with skew-symmetric torsion
 and special holonomy, \href{https://doi.org/10.1016/j.aim.2003.10.009}{\textit{Adv. Math.}} \textbf{189} (2004), 439--450,
 \href{https://arxiv.org/abs/math.DG/0302358}{arXiv:math.DG/0302358}.

\bibitem[9]{GJS}
Garcia-Fernandez M., Jordan J., Streets J., Non-{K}\"ahler {C}alabi--{Y}au
 geometry and pluriclosed flow, \href{https://arxiv.org/abs/2106.13716}{arXiv:2106.13716}.

\bibitem[10]{GS}
Garcia-Fernandez M., Streets J., Generalized {R}icci flow, \textit{Univ.
 Lecture Ser.}, Vol.~76, \href{https://doi.org/10.1090/ulect/076}{Amer. Math. Soc.}, Providence, RI, 2021.

\bibitem[11]{GHR}
Gates Jr. S.J., Hull C.M., Ro\v{c}ek M., Twisted multiplets and new
 supersymmetric nonlinear {$\sigma$}-models, \href{https://doi.org/10.1016/0550-3213(84)90592-3}{\textit{Nuclear Phys.~B}}
 \textbf{248} (1984), 157--186.

\bibitem[12]{Gau}
Gauduchon P., Hermitian connections and {D}irac operators, \textit{Boll. Un.
 Mat. Ital.~B} \textbf{11} (1997), 257--288.

\bibitem[13]{GGP}
Grantcharov D., Grantcharov G., Poon Y.S., Calabi--{Y}au connections with
 torsion on toric bundles, \href{https://doi.org/10.4310/jdg/1197320602}{\textit{J.~Differential Geom.}} \textbf{78} (2008),
 13--32, \href{https://arxiv.org/abs/math.DG/0306207}{arXiv:math.DG/0306207}.

\bibitem[14]{Grantcharov}
Grantcharov G., Geometry of compact complex homogeneous spaces with vanishing
 first {C}hern class, \href{https://doi.org/10.1016/j.aim.2010.10.005}{\textit{Adv. Math.}} \textbf{226} (2011), 3136--3159,
 \href{https://arxiv.org/abs/0905.0040}{arXiv:0905.0040}.

\bibitem[15]{GIP}
Gutowski J., Ivanov S., Papadopoulos G., Deformations of generalized
 calibrations and compact non-{K}\"ahler manifolds with vanishing first
 {C}hern class, \href{https://doi.org/10.4310/AJM.2003.v7.n1.a4}{\textit{Asian~J.~Math.}} \textbf{7} (2003), 39--79,
 \href{https://arxiv.org/abs/math.DG/0205012}{arXiv:math.DG/0205012}.

\bibitem[16]{HP}
Howe P.S., Papadopoulos G., Further remarks on the geometry of two-dimensional
 nonlinear {$\sigma$} models, \href{https://doi.org/10.1088/0264-9381/5/12/014}{\textit{Classical Quantum Gravity}} \textbf{5}
 (1988), 1647--1661.

\bibitem[17]{Hull}
Hull C.M., Compactifications of the heterotic superstring, \href{https://doi.org/10.1016/0370-2693(86)91393-6}{\textit{Phys.
 Lett.~B}} \textbf{178} (1986), 357--364.

\bibitem[24]{Samelson}
Samelson H., A class of complex-analytic manifolds, \textit{Portugal. Math.}
 \textbf{12} (1953), 129--132.

\bibitem[25]{Strominger}
Strominger A., Superstrings with torsion, \href{https://doi.org/10.1016/0550-3213(86)90286-5}{\textit{Nuclear Phys.~B}} \textbf{274}
 (1986), 253--284.

\bibitem[26]{Swann}
Swann A., Twisting {H}ermitian and hypercomplex geometries, \href{https://doi.org/10.1215/00127094-2010-059}{\textit{Duke
 Math.~J.}} \textbf{155} (2010), 403--431.

\bibitem[28]{Wang}
Wang H.-C., Closed manifolds with homogeneous complex structure,
 \href{https://doi.org/10.2307/2372397}{\textit{Amer.~J.~Math.}} \textbf{76} (1954), 1--32.

\end{thebibliography}
\end{document}
